% Part: first-order-logic
% Chapter: tableaux
% Section: rules-and-proofs

\documentclass[../../../include/open-logic-section]{subfiles}

\begin{document}

\iftag{FOL}
      {\olfileid{fol}{tab}{rul}}
      {\olfileid{pl}{tab}{rul}}

\olsection{قاعدې او \usetoken{P}{tableau}}

!!^a{tableau} په !!a{structure} کښې د~!!a{sentence} د رښتيا يا
دروغ کېدو د ممکنو لارو منظمه پلټنه ده۔ د تابلو بنسټيز توکي
!!{signed formula}s دي: !!{sentence}s له يوه رښتياوالي ارزښت
«نښې» سره، چې يا~$\True$ وي يا~$\False$۔ دا نښه لرونکي !!{formula}s
په داسې ونه کښې ترتيبېږي چې لاندې خوا ته غځېږي۔

\begin{defn}
  \emph{!!{signed formula}} له يوه رښتياوالي ارزښت او~!!a{sentence}
  څخه جوړه ده، يعنې له دې دواړو څخه يوه:
  \[
  \sFmla{\True}{!A} \text{ يا } \sFmla{\False}{!A}.
  \]
\end{defn}

په شهودي ډول $\sFmla{\True}{!A}$ داسې لوستلے شو چې «کېدے شي
$!A$ رښتيا وي»، او~$\sFmla{\False}{!A}$ داسې چې «کېدے شي
$!A$ دروغ وي» (په کوم !!{structure} کښې)۔

په ونه کښې هر !!{signed formula} يا \emph{فرض} دے (چې د ونې په
تر ټولو پاسنۍ برخه کښې ليکل کېږي)، يا تر ځان له پاسه له کوم
!!a{signed formula} څخه د يوې استنتاجي قاعدې په وسيله اخستل شوے
دے۔ د مخکښې !!{formula} د هر ممکن !!{main operator} دپاره دوه
قاعدې شته: يوه د هغه حالت دپاره چې نښه~$\True$ وي، او بله د هغه
حالت دپاره چې نښه~$\False$ وي۔ ځينې قاعدې ونې ته څانګې ورکوي او
ځينې يوازې څانګې ته !!{signed formula}s ورزياتوي۔ يوه قاعده
کېدے شي (او ډېر ځله بايد) سمدستي پر مخکښ !!{signed formula} نۀ،
بلکې له ريښې څخه د قاعدې د پلي کېدو تر ځاے پورې په څانګه کښې پر هر
!!{signed formula} پلې شي۔

يوه څانګه هله \emph{تړلې} ده چې
$\sFmla{\True}{!A}$ او~$\sFmla{\False}{!A}$ دواړه ولري۔
تړلے !!{tableau} هغه دے چې هره څانګه ئې تړلې وي۔ د شهودي تعبير له
مخې هره څانګه يو ګډ امکان بيانوي، خو
$\sFmla{\True}{!A}$ او~$\sFmla{\False}{!A}$ يو ځاے ممکن نۀ دي۔
په بله وينا، که يوه څانګه تړلې وي، هغه امکان رد شوے دے چې څانګه
ئې بيانوي۔ په ځانګړي ډول، د دې معنٰی دا ده چې تړلے !!{tableau} هغه
ټول امکانات ردوي چې د~$\sFmla{\True}{!A}$ په بڼه هر فرض په يو وخت
رښتيا او د~$\sFmla{\False}{!A}$ په بڼه هر فرض دروغ کړي۔

\emph{د~$!A$ دپاره} تړلے !!{tableau} داسې تړلے !!{tableau} دے چې
ريښه ئې~$\sFmla{\False}{!A}$ وي۔ که داسې تړلے !!{tableau} موجود
وي، د~$!A$ د دروغ کېدو ټول امکانات رد شوي؛ يعنې $!A$ بايد په هر
!!{structure} کښې رښتيا وي۔

\end{document}

